\input{../tex-inputs/latex-korrekturansicht-vorspann}

\section[Arthur Schnitzler an Theodor Herzl, 1. 7. 1893]{L03903 Arthur Schnitzler an Theodor Herzl, 1. 7. 1893}
\nopagebreak\mylabel{L03903v}
\rehead{ }\normalsize\beginnumbering\briefempfaengerindex{Herzl, Theodor@\textsc{Herzl, Theodor}!zzzSchnitzler, Arthur@\emph{von Arthur Schnitzler}!1893-07-011@{1. 7. 1893}|(be}
\toendnotes[C]{\smallbreak\pagebreak[2]}
\correspDesc{Versand  durch Arthur Schnitzler am 1. 7. 1893 in Wien
\newline{}Erhalt  durch Theodor Herzl im Zeitraum [2. 7. 1893 – 7. 7. 1893?] in Morschach}\toendnotes[C]{\smallbreak}
\Standort{Jerusalem, Central Zionist Archives, H1:1924-8.}
\physDesc{Brief, 1 Blatt, 2 Seiten, 304 Zeichen
\newline{}Handschrift: schwarze Tinte, deutsche Kurrent
\newline{}Ordnung: mit Bleistift von unbekannter Hand innerhalb das Konvoluts paginiert: »29« }
\pstart{}{\pb}Lieber Freund,\pend\vspace{0.5em}
\pstart
           ich theile Ihnen für alle Fälle mit, daſs ich heute auf 14 Tage nach \textcolor{pink}{Iſchl}\oindex{Bad Ischl@\textbf{Bad Ischl}|pw}{}\ledrightnote{\textcolor{pink}{Bad Ischl}} fahre, dort in der \textcolor{pink}{\textsc{Pension Leop. Petter}}\oindex{Hotel und Pension Rudolfshöhe (Leopold Petter)@\textbf{Hotel und Pension Rudolfshöhe (Leopold Petter)}, \emph{Hotel}|pw}{}\ledrightnote{\textcolor{pink}{Hotel und Pension Rudolfshöhe (Leopold Petter)}} wohne. Ich hoffe Sie {\pb}in \textcolor{pink}{Wien}\oindex{Wien@\textbf{Wien}, \emph{Verwaltungsgebiet}|pw}{}\ledrightnote{\textcolor{pink}{Wien}} oder in der Nähe jedenfalls noch anzutreffen. Haben sie doch die Güte,
               mir 2 Zeilen nach \textcolor{pink}{Iſchl}\oindex{Bad Ischl@\textbf{Bad Ischl}|pw}{}\ledrightnote{\textcolor{pink}{Bad Ischl}}. zu ſchreiben.\pend
           
\pstart
           Mit herzlich Grüßen{\\[\baselineskip]}Ihr \spacefill\mbox{ArthSchnitzler}\pend
           \leftskip=0em{}
\pstart
           1. 7. 93.\pend
           \selectlanguage{ngerman}\endnumbering\briefempfaengerindex{Herzl, Theodor@\textsc{Herzl, Theodor}!zzzSchnitzler, Arthur@\emph{von Arthur Schnitzler}!1893-07-011@{1. 7. 1893}|)be}\mylabel{L03903h}  \newcommand{\dateiname}{L03903}\newcommand{\titel}{Arthur Schnitzler an Theodor Herzl, 1. 7. 1893}\newcommand{\editorInnen}{Herausgegeben von Jahnke, SelmaMüller, Martin Anton}\input{../tex-inputs/latex-korrekturansicht-abspann}
